\chapter{Integrability And The Rigid Body}
\label{chap:rigid-integrability}

\section{Why The Rigid Body Is Central}
\label{sec:rigid-why-central}

The free rigid body is a nearly ideal teaching example. It is nonlinear, it
lives naturally on the Lie group \(\SO(3)\), it has a reduced Hamiltonian
description on angular momentum space, and it is completely integrable. At the
same time, small modifications such as gravity acting on a heavy top lead toward
nonintegrability and chaos.

\begin{assumption}[Rigid body model]
A rigid body is an ideal continuum whose internal distances do not change. A
material point is labelled by a body coordinate \(X\in B\subset\R^3\). Its
spatial position is
\[
  x(t)=r(t)+R(t)X,
\]
where \(r(t)\in\R^3\) is the center-of-mass position and
\(R(t)\in\SO(3)\) is the attitude matrix. The free rigid body chapter suppresses
translation by working in the center-of-mass frame.
\end{assumption}

\begin{table}[htbp]
\centering
\caption{Rigid-body and integrability notation. Symbols that are overloaded in
other chapters are cross-listed in Section~\ref{sec:overloaded-symbols}.}
\label{tab:rigid-integrability-notation}
\small
\begin{tabular}{p{0.18\textwidth}p{0.48\textwidth}p{0.25\textwidth}}
\toprule
Symbol & Meaning & Space or units\\
\midrule
\(R\) & attitude matrix & \(\SO(3)\)\\
\(\hat e_a\) & \(a\)-th body axis seen in space & unit vector in \(\R^3\)\\
\(\omega_{\mathrm{space}}\) & spatial angular velocity & \(\R^3\), \(T^{-1}\)\\
\(\Omega\) & body angular velocity & \(\R^3\), \(T^{-1}\)\\
\(\widehat\Omega\) & skew matrix with \(\widehat\Omega v=\Omega\times v\) &
\(\mathfrak{so}(3)\)\\
\(I\) & inertia tensor in body frame & symmetric positive matrix\\
\(I_1,I_2,I_3\) & principal moments of inertia & positive scalars\\
\(M\) & body angular momentum & \(I\Omega\)\\
\(J\) & spatial angular momentum & \(J=RM\)\\
\(H\) & kinetic energy Hamiltonian & scalar energy\\
\(C\) & squared angular momentum & Casimir on \(\mathfrak{so}(3)^*\)\\
\(\ell(\Omega)\) & reduced Lagrangian & function on \(\mathfrak{so}(3)\)\\
\(F_i\) & commuting first integrals & Liouville theorem\\
\(M_f\) & common regular level set \(F_i=f_i\) & invariant manifold\\
\(\Lambda\) & period lattice of commuting flows & subgroup of \(\R^n\)\\
\(\gamma_j\) & one-cycle on an invariant torus & element of \(H_1(\mathbb T^n)\)\\
\(I_j\) & action variable & \((2\pi)^{-1}\oint_{\gamma_j}p_i\dd q^i\)\\
\(\theta^j\) & angle variable conjugate to \(I_j\) & \(2\pi\)-periodic coordinate\\
\(S(I,q)\) & generating function for angles & locally multivalued\\
\bottomrule
\end{tabular}
\end{table}

\paragraph{Roadmap.}
The chapter moves from a concrete body to the general geometry of integrable
systems. The first half explains the kinematics and dynamics of a free rigid
body: rotations live on \(\SO(3)\), angular velocity lives in its Lie algebra,
and the reduced motion lives on angular momentum spheres. The second half
abstracts the lesson: enough independent commuting first integrals force
regular compact level sets to be tori, and on those tori the motion is linear.
The rigid body is therefore not just an example after the theorem; it is the
model that makes the theorem visible.

\paragraph{Three layers to keep separate.}
Rigid-body mechanics mixes three descriptions that are often compressed into
one line of notation. The kinematic layer says how \(R(t)\in\SO(3)\) defines
angular velocity. The dynamical layer says how the kinetic energy and the
inertia tensor produce Euler's equations. The reduced Hamiltonian layer says
why angular momentum spheres carry a Lie-Poisson flow. Each layer has its own
variables and its own conserved quantities. Confusing them is the fastest way
to lose a sign or to mistake body components for spatial components.

\section{The Rotation Group And Angular Velocity}
\label{sec:rotation-group-angular-velocity}

The rotation group is
\[
  \SO(3)=\{R\in\R^{3\times 3}:R^TR=I,\ \det R=1\}.
\]
It is useful to think of \(R\) as the matrix whose columns are the three
orthonormal body axes \(\hat e_1,\hat e_2,\hat e_3\) as seen in space:
\[
  R=(\hat e_1\ \hat e_2\ \hat e_3).
\]
Thus a point of configuration space is not a triple of unconstrained vectors;
it is an orthonormal oriented frame. The condition \(R^TR=I\) gives six
constraints on nine matrix entries, leaving three degrees of freedom.

The spatial angular velocity \(\omega_{\mathrm{space}}\) is defined by the fact that every body
axis instantaneously rotates according to
\[
  \frac{\dd \hat e_a}{\dd t}=\omega_{\mathrm{space}}\times\hat e_a,
\]
or, in components,
\[
  \dot R_{ia}=\epsilon_{ijk}(\omega_{\mathrm{space}})_jR_{ka}.
\]

Differentiate \(R^TR=I\):
\[
  \dot R^TR+R^T\dot R=0.
\]
Thus \(R^T\dot R\) is skew-symmetric. Every skew-symmetric \(3\times3\) matrix
is uniquely of the form
\[
  \widehat\Omega=
  \begin{pmatrix}
  0 & -\Omega_3 & \Omega_2\\
  \Omega_3 & 0 & -\Omega_1\\
  -\Omega_2 & \Omega_1 & 0
  \end{pmatrix},
\]
where \(\widehat\Omega v=\Omega\times v\). The body angular velocity is defined
by
\[
  R^T\dot R=\widehat\Omega.
\]
The spatial angular velocity \(\omega_{\mathrm{space}}\) satisfies
\[
  \dot R R^T=\widehat{\omega}_{\mathrm{space}},
  \qquad
  \omega_{\mathrm{space}}=R\Omega.
\]

This space/body distinction is one of the main sources of sign errors in rigid
body mechanics. The vector \(\omega_{\mathrm{space}}\) is expressed in fixed
space axes, while \(\Omega\) is expressed in the moving body axes. They describe
the same instantaneous rotation, but derivatives of body components must also
differentiate the moving frame. Euler's equation is exactly where this
bookkeeping becomes dynamical.

\section{\texorpdfstring{Euler Angles, Metric, And The Topology Of \(\SO(3)\)}{Euler Angles, Metric, And The Topology Of SO(3)}}
\label{sec:euler-angles-so3}

One can parametrize a dense open set of \(\SO(3)\) by Euler angles
\[
  (\phi,\theta,\psi),
  \qquad
  0\leq \phi<2\pi,\quad 0\leq \theta\leq \pi,\quad 0\leq \psi<2\pi.
\]
For the Euler-angle convention used below, the body axes are
\[
\begin{aligned}
\hat e_1 &=
\begin{pmatrix}
\cos\phi\cos\psi-\sin\phi\cos\theta\sin\psi\\
\sin\phi\cos\psi+\cos\phi\cos\theta\sin\psi\\
\sin\theta\sin\psi
\end{pmatrix},\\
\hat e_2 &=
\begin{pmatrix}
-\cos\phi\sin\psi-\sin\phi\cos\theta\cos\psi\\
-\sin\phi\sin\psi+\cos\phi\cos\theta\cos\psi\\
\sin\theta\cos\psi
\end{pmatrix},\\
\hat e_3 &=
\begin{pmatrix}
\sin\phi\sin\theta\\
-\cos\phi\sin\theta\\
\cos\theta
\end{pmatrix}.
\end{aligned}
\]
The exact convention is less important than the lesson: Euler angles are a
coordinate chart, not the space itself. The chart is singular at
\(\theta=0,\pi\), even though \(\SO(3)\) is a smooth homogeneous space.

A coordinate-independent line element on \(\SO(3)\) is
\[
  \dd s^2=-\frac12\tr\left(R^{-1}\dd R\right)^2.
\]
For this standard Euler-angle convention, this becomes
\[
  \dd s^2
  =
  \dd\phi^2+\dd\theta^2+\dd\psi^2
  +2\cos\theta\,\dd\phi\,\dd\psi.
\]
Introduce
\[
  \xi_1=\phi+\psi,\qquad \xi_2=\phi-\psi.
\]
Then
\[
  \dd s^2
  =
  \dd\theta^2
  +
  \cos^2\frac{\theta}{2}\,\dd\xi_1^2
  +
  \sin^2\frac{\theta}{2}\,\dd\xi_2^2.
\]
This formula shows the topology pictorially. At one endpoint of the
\(\theta\)-interval the \(\xi_2\)-circle shrinks; at the other endpoint the
\(\xi_1\)-circle shrinks. Gluing these solid tori gives \(S^3\), and the
remaining identification is the antipodal identification. Hence
\[
  \SO(3)\simeq S^3/\Z_2.
\]
This is why a rigid body admits local angle coordinates but no single global
nondegenerate Euler-angle chart.

\section{Kinetic Energy And The Inertia Tensor}
\label{sec:inertia-tensor}

Take the center of mass as origin in body coordinates:
\[
  \int_B X\,\rho(X)\,\dd^3X=0,
\]
where \(\rho(X)\) is the mass density. A point \(X\) has velocity
\[
  \dot x=\dot R X=R\widehat\Omega X=R(\Omega\times X).
\]
Since \(R\) preserves Euclidean lengths,
\[
  \norm{\dot x}^2=\norm{\Omega\times X}^2.
\]
The kinetic energy is
\[
  T=\frac12\int_B\rho(X)\norm{\Omega\times X}^2\,\dd^3X.
\]
Using
\[
  \norm{\Omega\times X}^2
  =
  \norm{\Omega}^2\norm{X}^2-(\Omega\cdot X)^2,
\]
we obtain
\[
  T=\frac12\Omega^T I\Omega,
\]
where
\[
  I_{ij}
  =
  \int_B\rho(X)
  \left(\norm{X}^2\delta_{ij}-X_iX_j\right)\dd^3X.
\]
The matrix \(I\) is the inertia tensor in the body frame. In principal axes,
\[
  I=\diag(I_1,I_2,I_3),
\]
and
\[
  T=\frac12(I_1\Omega_1^2+I_2\Omega_2^2+I_3\Omega_3^2).
\]
Equivalently, using the spatial angular velocity \(\omega_{\mathrm{space}}\) and the body axes
\(\hat e_a\),
\[
  L=\frac12\sum_{a=1}^3 I_a(\hat e_a\cdot\omega_{\mathrm{space}})^2.
\]
The spatial angular momentum is
\[
  J_i=\frac{\partial L}{\partial (\omega_{\mathrm{space}})_i}
  =
  \sum_{a=1}^3 I_a(\hat e_a)_i(\hat e_a\cdot\omega_{\mathrm{space}}).
\]
Its body components are
\[
  \widetilde J_a=\hat e_a\cdot J=I_a\Omega_a,
\]
and the Hamiltonian is
\[
  H=\sum_{a=1}^3\frac{\widetilde J_a^2}{2I_a}.
\]

\section{\texorpdfstring{Full Phase Space: \(T^*\SO(3)\)}{Full Phase Space: T*SO(3)}}
\label{sec:rigid-full-phase-space}

The Hamiltonian phase space of a finite-dimensional mechanical system with
configuration space \(Q\) is the cotangent bundle \(T^*Q\). Therefore the
unreduced phase space of a rigid body with fixed center of mass is
\[
  T^*\SO(3).
\]
If \(q^i\), \(i=1,2,3\), are any nondegenerate local coordinates on
\(\SO(3)\), the canonical momenta are
\[
  p_i=\frac{\partial L}{\partial \dot q^i}.
\]
Under a change of coordinates \(q^i=q^i(\tilde q)\), the momenta transform as a
one-form:
\[
  \tilde p_j
  =
  \frac{\partial q^i}{\partial \tilde q^j}p_i.
\]
This is the coordinate calculation behind the intrinsic statement
\((q,p)\in T^*Q\).

\begin{remark}[Angular momentum is not a canonical momentum]
Let \(M_a\) be the components of body angular momentum in principal axes:
\[
  M_a=I_a\Omega_a.
\]
The \(M_a\) are excellent coordinates on the fiber of \(T^*\SO(3)\) after a
choice of body frame, but they are not canonical momenta conjugate to Euler
angles. Locally one has a nondegenerate linear relation
\[
  p_i=A_{ia}(q)M_a,
\]
where the matrix \(A(q)\) depends on the chosen angle coordinates. The canonical
symplectic form is still
\[
  \omega=\dd q^i\wedge \dd p_i,
\]
but when it is expressed in the noncanonical variables \(M_a\), the Poisson
brackets become the rigid-body Lie-Poisson brackets rather than
\(\Poisson{M_a}{M_b}=0\).
Globally, because \(\SO(3)\) is a Lie group, left or right translation
trivializes the cotangent bundle:
\[
  T^*\SO(3)\simeq \SO(3)\times\mathfrak{so}(3)^*.
\]
This global trivialization is the precise sense in which the angular-momentum
components give global fiber coordinates, even though they are noncanonical.
\end{remark}

\section{\texorpdfstring{Variations On \(\SO(3)\)}{Variations On SO(3)}}
\label{sec:variations-so3}

A variation of \(R(t)\in\SO(3)\) must remain tangent to the group. Write
\[
  \delta R=R\widehat\Sigma,
\]
where \(\Sigma(t)\in\R^3\) vanishes at the time endpoints. The induced variation
of body angular velocity is
\[
  \delta\widehat\Omega
  =
  \delta(R^T\dot R)
  =
  \dot{\widehat\Sigma}+[\widehat\Omega,\widehat\Sigma].
\]
Under the hat isomorphism this becomes
\[
  \delta\Omega=\dot\Sigma+\Omega\times\Sigma.
\]

\begin{derivation}[Euler-Poincare equation for the free rigid body]
The reduced action is
\[
  S[R]=\int_{t_0}^{t_1}\ell(\Omega)\,\dd t,
  \qquad
  \ell(\Omega)=\frac12\Omega^TI\Omega.
\]
Let
\[
  M=\frac{\partial \ell}{\partial \Omega}=I\Omega.
\]
Then
\[
\begin{aligned}
  \delta S
  &=
  \int_{t_0}^{t_1} M\cdot(\dot\Sigma+\Omega\times\Sigma)\,\dd t\\
  &=
  \left[M\cdot\Sigma\right]_{t_0}^{t_1}
  -
  \int_{t_0}^{t_1}\dot M\cdot\Sigma\,\dd t
  +
  \int_{t_0}^{t_1}M\cdot(\Omega\times\Sigma)\,\dd t.
\end{aligned}
\]
The boundary term vanishes. Since
\[
  M\cdot(\Omega\times\Sigma)=(M\times\Omega)\cdot\Sigma,
\]
stationarity for arbitrary \(\Sigma\) gives
\[
  -\dot M+M\times\Omega=0.
\]
Equivalently,
\[
  \boxed{\dot M+\Omega\times M=0}
\]
with the hat convention \(\widehat\Omega v=\Omega\times v\) used throughout
these notes. In components this is
\[
  I_1\dot\Omega_1=(I_2-I_3)\Omega_2\Omega_3,
\]
and cyclic permutations.
\end{derivation}

\section{Euler Equations In Principal Axes}
\label{sec:euler-equations-principal}

With \(M_i=I_i\Omega_i\), Euler's equations are
\[
\begin{aligned}
  I_1\dot\Omega_1 &= (I_2-I_3)\Omega_2\Omega_3,\\
  I_2\dot\Omega_2 &= (I_3-I_1)\Omega_3\Omega_1,\\
  I_3\dot\Omega_3 &= (I_1-I_2)\Omega_1\Omega_2.
\end{aligned}
\]
These equations are autonomous and quadratic. They conserve the kinetic energy
\[
  H=\frac12(I_1\Omega_1^2+I_2\Omega_2^2+I_3\Omega_3^2)
  =
  \frac12\left(\frac{M_1^2}{I_1}+\frac{M_2^2}{I_2}+\frac{M_3^2}{I_3}\right)
\]
and the squared angular momentum
\[
  C=\norm{M}^2=M_1^2+M_2^2+M_3^2.
\]

\begin{proof}[Direct conservation check]
Since \(M=I\Omega\),
\[
  \dot H=\Omega\cdot \dot M.
\]
Using \(\dot M=-\Omega\times M\),
\[
  \dot H=-\Omega\cdot(\Omega\times M)=0.
\]
Similarly,
\[
  \dot C=2M\cdot\dot M=-2M\cdot(\Omega\times M)=0.
\]
\end{proof}

\section{Geometry Of The Reduced Motion}
\label{sec:rigid-reduced-geometry}

The reduced motion in \(M\)-space lies on the intersection of:
\[
  C(M)=\norm{M}^2=\text{constant}
\]
and
\[
  H(M)=\frac12M^TI^{-1}M=\text{constant}.
\]
Thus it is the intersection of an angular momentum sphere and an energy
ellipsoid.

\begin{figure}[htbp]
\centering
\begin{tikzpicture}[
    scale=1.05,
    >=Latex,
    x={(1cm,0cm)},
    y={(0.42cm,0.27cm)},
    z={(0cm,1cm)},
    line join=round
  ]
  \draw[->, gray!65] (0,0,0) -- (2.75,0,0) node[right] {\(M_1\)};
  \draw[->, gray!65] (0,0,0) -- (0,2.45,0) node[above right] {\(M_2\)};
  \draw[->, gray!65] (0,0,0) -- (0,0,3.15) node[above] {\(M_3\)};

  \draw[red!65!black, thick, domain=0:360, samples=120, variable=\t]
    plot ({2*cos(\t)}, {2*sin(\t)}, {0});
  \draw[red!65!black, dashed, domain=0:360, samples=120, variable=\t]
    plot ({2*cos(\t)}, {0}, {2*sin(\t)});
  \draw[red!65!black, dashed, domain=0:360, samples=120, variable=\t]
    plot ({0}, {2*cos(\t)}, {2*sin(\t)});

  \draw[blue!70!black, thick, domain=0:360, samples=140, variable=\t]
    plot ({sqrt(847/340)*cos(\t)}, {sqrt(2541/507)*sin(\t)}, {0});
  \draw[blue!70!black, thin, domain=0:360, samples=140, variable=\t]
    plot ({sqrt(847/340)*cos(\t)}, {0}, {2.75*sin(\t)});
  \draw[blue!70!black, thin, domain=0:360, samples=140, variable=\t]
    plot ({0}, {sqrt(2541/507)*cos(\t)}, {2.75*sin(\t)});

  \draw[purple!80!black, very thick, domain=0:360, samples=180, variable=\t]
    plot ({sqrt(1 + 0.75*sin(\t)*sin(\t))}, {1.732*cos(\t)}, {1.5*sin(\t)});
  \draw[purple!80!black, very thick, dashed, domain=0:360, samples=180, variable=\t]
    plot ({-sqrt(1 + 0.75*sin(\t)*sin(\t))}, {1.732*cos(\t)}, {1.5*sin(\t)});

  \fill[purple!80!black] ({sqrt(1 + 0.75*0*0)}, {1.732}, {0}) circle (1.3pt);
  \node[red!65!black, align=center] at (-2.2,-1.15,-0.65)
    {angular momentum\\sphere \(C=\norm M^2\)};
  \node[blue!70!black, align=center] at (1.95,1.55,2.95)
    {energy ellipsoid\\\(2H=M^TI^{-1}M\)};
  \node[purple!80!black] at (-0.55,-2.2,1.45)
    {intersection trajectories};
\end{tikzpicture}
\caption{Rigid-body reduction in angular-momentum space. Conservation of
\(\norm M^2\) confines the motion to a sphere, conservation of \(H\) confines it
to an ellipsoid, and the reduced trajectory is their intersection. The solid
and dashed purple curves are the two projected branches of the exact
three-dimensional intersection for the sphere \(\norm M^2=4\) and the
representative ellipsoid with semi-axes
\(\sqrt{847/340},\sqrt{2541/507},11/4\).}
\label{fig:rigid-body-reduced-geometry}
\end{figure}

\section{Stability Of Principal-Axis Rotations}
\label{sec:principal-axis-stability}

Assume
\[
  I_1<I_2<I_3.
\]
A steady rotation about a principal axis is a solution. For example,
\[
  \Omega(t)=(\Omega_0,0,0)
\]
is rotation about the first axis.

\begin{derivation}[Linear stability]
Perturb rotation about the first axis:
\[
  \Omega_1=\Omega_0+\alpha_1,\qquad
  \Omega_2=\alpha_2,\qquad
  \Omega_3=\alpha_3,
\]
and keep only linear terms. Then
\[
  \dot\alpha_1=0,
\]
while
\[
  I_2\dot\alpha_2=(I_3-I_1)\Omega_0\alpha_3,
  \qquad
  I_3\dot\alpha_3=(I_1-I_2)\Omega_0\alpha_2.
\]
Differentiate the first of these:
\[
  \ddot\alpha_2
  =
  \frac{(I_3-I_1)(I_1-I_2)}{I_2I_3}\Omega_0^2\alpha_2.
\]
Since \(I_3-I_1>0\) and \(I_1-I_2<0\), the coefficient is negative, so
\(\alpha_2\) oscillates. Rotation about the first axis is linearly stable.

For rotation about the second axis, the analogous coefficient is
\[
  \frac{(I_2-I_3)(I_1-I_2)}{I_1I_3}\Omega_0^2>0,
\]
because both factors in the numerator are negative. The perturbation therefore
has exponential solutions. Thus the intermediate axis is unstable. Rotation
about the third axis is stable.
\end{derivation}

The geometric reason is that the energy restricted to a fixed angular momentum
sphere has extrema at the smallest and largest inertia axes but a saddle at the
intermediate axis.

\section{Lie-Poisson Structure}
\label{sec:lie-poisson-rigid-body}

The reduced rigid body is Hamiltonian on \(\mathfrak{so}(3)^*\cong\R^3\), but
not with canonical coordinates. The Lie-Poisson bracket is
\[
  \Poisson{F}{G}(M)
  =
  -M\cdot\left(\nabla F(M)\times\nabla G(M)\right).
\]
For
\[
  H(M)=\frac12M\cdot I^{-1}M,
  \qquad
  \nabla H=I^{-1}M=\Omega,
\]
the equation for any observable \(F\) is
\[
  \dot F=\Poisson{F}{H}.
\]
Taking \(F(M)=M_i\) yields
\[
  \dot M=-\Omega\times M.
\]

\begin{definition}[Casimir]
A Casimir is a function \(C\) satisfying
\[
  \Poisson{C}{F}=0
\]
for every observable \(F\). For the rigid body bracket,
\[
  C(M)=\norm{M}^2
\]
is a Casimir. Its level sets are the coadjoint orbits, which here are spheres.
\end{definition}

\section{Liouville Integrability And Invariant Tori}
\label{sec:liouville-tori}

The rigid-body calculations above show a pattern: conservation laws confine
motion to lower-dimensional sets, and the reduced flow on those sets can be
simple even when the original coordinates look complicated. Liouville
integrability is the precise Hamiltonian version of this pattern. It is not
just ``having many constants of motion.'' The constants must be independent,
they must Poisson commute, and their common level set must be regular enough to
carry the commuting Hamiltonian flows.

The three hypotheses have separate meanings. Independence says the level set
has the expected dimension. Involution says the Hamiltonian flows generated by
the integrals are compatible rather than fighting one another. Compactness says
the flow coordinates eventually repeat, turning straight-line motion on
\(\R^n\) into straight-line motion on a torus. Dropping any one of these
assumptions changes the conclusion.

\begin{definition}[Liouville integrability]
A Hamiltonian system on a \(2n\)-dimensional symplectic manifold is Liouville
integrable if there are \(n\) independent functions
\[
  F_1=H,\ F_2,\ldots,F_n
\]
such that
\[
  \Poisson{F_i}{F_j}=0
\]
for all \(i,j\). The functions are said to be in involution.
\end{definition}

\begin{example}[Basic integrable systems]
Any autonomous Hamiltonian system with one degree of freedom is integrable: the
Hamiltonian itself supplies the one conserved quantity. The fixed-COM rigid body
has three degrees of freedom. The three spatial functions
\[
  H,\qquad J_3,\qquad J^2
\]
are independent and mutually Poisson commuting on regular sets, even though the
individual components \(J_1,J_2,J_3\) do not commute with one another. This is
the rigid-body example used throughout this chapter.
\end{example}

\begin{assumption}[Regular compact level set]
Fix constants \(f=(f_1,\ldots,f_n)\) and define
\[
  M_f=\{x\in P: F_i(x)=f_i,\ i=1,\ldots,n\},
\]
where \(P\) is the \(2n\)-dimensional phase space. We assume that the
differentials \(\dd F_i\) are independent on \(M_f\) and that \(M_f\) is
compact and connected.
\end{assumption}

The regularity assumption says that \(M_f\) is a smooth \(n\)-dimensional
submanifold. Its tangent vectors \(v\in T_xM_f\) are exactly those satisfying
\[
  \dd F_i(v)=0
  \qquad\text{for all }i.
\]
Let \(X_{F_i}\) be the Hamiltonian vector field of \(F_i\), defined by
\[
  \iota_{X_{F_i}}\omega=\dd F_i.
\]
Then
\[
  \dd F_j(X_{F_i})
  =
  \Poisson{F_j}{F_i}
  =
  0.
\]
Thus every \(X_{F_i}\) is tangent to \(M_f\). The commutator of two Hamiltonian
vector fields is
\[
  [X_{F_i},X_{F_j}]
  =
  -X_{\Poisson{F_i}{F_j}},
\]
so involution implies
\[
  [X_{F_i},X_{F_j}]=0.
\]
We therefore get \(n\) independent commuting vector fields on \(M_f\).

\begin{derivation}[Why the regular compact level set is a torus]
The proof has a simple geometric shape. The Hamiltonian vector fields
\(X_{F_i}\) give \(n\) independent directions tangent to \(M_f\), and because
the integrals Poisson commute, flowing in one direction and then another gives
the same point as flowing in the opposite order. Thus \(M_f\) is locally
coordinatized by ``flow times.'' Compactness is what turns these flow-time
coordinates from \(\R^n\) into periodic coordinates.

Commuting vector fields have commuting flows. Starting at a point \(x_0\in
M_f\), flow for times \(s=(s^1,\ldots,s^n)\) along the \(n\) vector fields:
\[
  \Phi_s(x_0)
  =
  \exp(s^1X_{F_1})\cdots\exp(s^nX_{F_n})x_0.
\]
Because the flows commute, this defines an action of \(\R^n\) on \(M_f\).
Independence of the vector fields makes the action locally free, so the
coordinates \(s^i\) are local coordinates on \(M_f\). The only possible global
identifications are translations
\[
  s\sim s+\lambda,\qquad \lambda\in\Lambda,
\]
where \(\Lambda\subset\R^n\) is the period lattice of the orbit. Compactness
forces \(\Lambda\) to have rank \(n\): if the rank were smaller, the quotient
\(\R^n/\Lambda\) would still contain at least one noncompact Euclidean
direction. The action is locally free, so each orbit is open in the
\(n\)-manifold \(M_f\). Since \(M_f\) is connected, that open orbit is the
whole level set. Hence
\[
  M_f\simeq \R^n/\Lambda\simeq \mathbb T^n.
\]
These are the invariant Liouville tori.
\end{derivation}

This argument is often the first place where the word ``torus'' becomes more
than a picture. The torus is not assumed; it is forced by the existence of
enough commuting flows on a compact regular level set. The angle variables
introduced below are coordinates along these flows, and the action variables
are the quantities that label which torus one is on.

\begin{figure}[htbp]
\centering
\begin{tikzpicture}[scale=1.0, >=Latex]
  \draw[thick] (0,0) rectangle (4.2,3.0);
  \draw[->, blue!70!black, thick] (0.45,0.55) -- (3.75,0.55)
    node[midway, below] {\(\gamma_1\)};
  \draw[->, red!75!black, thick] (0.55,0.45) -- (0.55,2.55)
    node[midway, left] {\(\gamma_2\)};
  \draw[purple!75!black, thick, ->]
    (0.25,0.25) -- (1.1,0.9) -- (1.95,1.55) -- (2.8,2.2) -- (3.65,2.85);
  \draw[purple!75!black, thick, ->]
    (3.65,2.85) -- (4.15,3.23);
  \draw[purple!75!black, thick, ->]
    (-0.15,-0.23) -- (0.25,0.25);
  \draw[gray!70, ->] (4.2,1.5) -- (4.8,1.5)
    node[right] {identify};
  \draw[gray!70, ->] (2.1,3.0) -- (2.1,3.55)
    node[above] {identify};
  \node[align=center] at (2.1,-0.75)
    {flat model of a two-dimensional invariant torus\\
     \(\mathbb R^2/\Lambda\) with angle coordinates};
  \node[purple!75!black, align=center] at (5.25,2.45)
    {Hamiltonian flow\\straightens to\\constant slope};
\end{tikzpicture}
\caption{A Liouville torus is locally a space of commuting flow times with
periodic identifications. Action variables label the torus; angle variables
are coordinates along the fundamental cycles.}
\label{fig:liouville-flat-torus}
\end{figure}

The symplectic form vanishes when restricted to an invariant torus. Indeed, on
the tangent basis \(X_{F_i}\),
\[
  \omega(X_{F_i},X_{F_j})
  =
  \dd F_i(X_{F_j})
  =
  \Poisson{F_i}{F_j}
  =
  0.
\]
Since these vectors span \(T_xM_f\), the torus is Lagrangian:
\[
  \omega|_{M_f}=0.
\]

\begin{derivation}[Coordinate proof that the invariant torus is Lagrangian]
A useful coordinate check is the following. Use canonical
coordinates \((q_i,p_i)\), and suppose the level set is written locally as
\[
  F_i(p,q)=f_i,\qquad i=1,\ldots,n.
\]
If the equations can be solved locally for
\[
  p_i=p_i(q,F),
\]
then
\[
\begin{aligned}
  \omega
  &=
  \dd q_i\wedge\dd p_i\\
  &=
  \dd q_i\wedge
  \left(
    \frac{\partial p_i}{\partial q_j}\bigg|_F\dd q_j
    +
    \frac{\partial p_i}{\partial F_j}\bigg|_q\dd F_j
  \right).
\end{aligned}
\]
The key claim is
\[
  \frac{\partial p_i}{\partial q_j}\bigg|_F
  =
  \frac{\partial p_j}{\partial q_i}\bigg|_F.
\]
Then the \(\dd q_i\wedge\dd q_j\) contribution vanishes by antisymmetry, and
\[
  \omega=
  \frac{\partial p_i}{\partial F_j}\bigg|_q\dd q_i\wedge\dd F_j.
\]
On the invariant torus \(F_j=f_j\), \(\dd F_j=0\), so
\[
  \omega|_{M_f}=0.
\]

To prove the claim, differentiate the identities \(F_k(p(F,q),q)=F_k\) at
fixed \(F\):
\[
  0=
  \frac{\partial F_k}{\partial p_j}\bigg|_q
  \frac{\partial p_j}{\partial q_i}\bigg|_F
  +
  \frac{\partial F_k}{\partial q_i}\bigg|_p.
\]
Let
\[
  A_{kj}=\frac{\partial F_k}{\partial p_j}\bigg|_q,\qquad
  M_{ji}=\frac{\partial p_j}{\partial q_i}\bigg|_F,\qquad
  B_{ki}=\frac{\partial F_k}{\partial q_i}\bigg|_p.
\]
Then
\[
  AM+B=0,\qquad M=-A^{-1}B.
\]
The involution condition gives
\[
  0=\Poisson{F_i}{F_j}
  =
  A_{ik}B_{jk}-A_{jk}B_{ik},
\]
or
\[
  AB^T-BA^T=0.
\]
Hence
\[
  A^{-1}B=B^T(A^T)^{-1}=(A^{-1}B)^T,
\]
so \(M=M^T\), proving the symmetry claim.
\end{derivation}

\section{Action Variables And Angle Variables}
\label{sec:action-angle-variables}

Choose a basis of one-cycles \(\gamma_1,\ldots,\gamma_n\) on an invariant torus
\(M_f\). In canonical coordinates the action variables are
\[
  I_j=\frac{1}{2\pi}\oint_{\gamma_j} p_i\,\dd q^i.
\]
Actions are therefore not arbitrary constants of integration. They are
phase-space areas measured around the independent cycles of the torus. This is
why they are insensitive to small deformations of the cycle and why they
survive as adiabatic invariants in Chapter~\ref{chap:hamilton-jacobi}.
They are invariant under continuous deformations of \(\gamma_j\) on the torus.
If \(\gamma_j\) and \(\gamma_j'\) bound a two-chain \(D\subset M_f\), then by
Stokes' theorem
\[
  \oint_{\gamma_j}p_i\,\dd q^i
  -
  \oint_{\gamma_j'}p_i\,\dd q^i
  =
  \int_D \dd(p_i\,\dd q^i)
  =
  \int_D \omega
  =
  0.
\]
Thus \(I_j\) depends on the torus, not on the representative loop. Under
regularity assumptions, the \(I_j\) can replace the original integrals \(F_j\)
as labels for nearby tori.

To find conjugate angles, construct a local generating function
\[
  S(I,q)
  =
  \int^q p_i(I,q')\,\dd q'^i,
\]
where the integration is performed on a local branch. The compatibility
condition for this definition is
\[
  \frac{\partial p_i}{\partial q^j}\bigg|_I
  =
  \frac{\partial p_j}{\partial q^i}\bigg|_I,
\]
which is another form of \(\omega|_{M_I}=0\). Define
\[
  \theta^j=\frac{\partial S}{\partial I_j}\bigg|_q.
\]
When one goes around the cycle \(\gamma_k\), the multivalued function \(S\)
jumps by
\[
  \Delta_{\gamma_k}S=\oint_{\gamma_k}p_i\,\dd q^i=2\pi I_k.
\]
Therefore
\[
  \Delta_{\gamma_k}\theta^j
  =
  \frac{\partial}{\partial I_j}(2\pi I_k)
  =
  2\pi\delta^j{}_k.
\]
The \(\theta^j\) are angles on the torus:
\[
  \theta^j\sim \theta^j+2\pi.
\]
In these action-angle coordinates the Hamiltonian is a function of the actions
alone,
\[
  H=H(I),
\]
and Hamilton's equations reduce to
\[
  \dot I_j=0,\qquad
  \dot\theta^j=\omega^j(I)=\frac{\partial H}{\partial I_j}.
\]
Thus motion on each invariant torus is linear. If the frequencies are rationally
related, the orbit closes; otherwise the orbit is dense on the torus.

\section{Rigid Body Tori}
\label{sec:rigid-body-tori}

The full rigid body with fixed center of mass has three degrees of freedom. On
\(T^*\SO(3)\) the functions
\[
  H,\qquad J_3,\qquad J^2
\]
are mutually Poisson commuting when \(J_3\) is a fixed spatial component of
angular momentum and \(J^2=\norm{J}^2\). For regular values, their common level
set is a three-torus.

The geometry can be seen in two stages. First, the body angular momentum
\(M=R^TJ\) is constrained by
\[
  J^2=M_1^2+M_2^2+M_3^2=\text{constant}
\]
and
\[
  H=\frac12\left(\frac{M_1^2}{I_1}
  +\frac{M_2^2}{I_2}
  +\frac{M_3^2}{I_3}\right)=\text{constant}.
\]
This puts \(M\) on the closed curve obtained by intersecting a sphere with an
ellipsoid. Second, fixing \(J_3\) constrains the orientation \(R\): the third
spatial component of \(J=RM\) is fixed. For a given \(M\) on the curve, this
condition leaves a circle of allowed attitudes. A circle fibered over the
momentum-space circle gives a two-dimensional torus in the reduced picture, and
including the remaining cyclic angle gives the full three-torus.

The reduced free rigid body has one effective degree of freedom on each angular
momentum sphere. The energy level on that sphere gives closed curves except at
separatrices. This is the concrete meaning of integrability here.

\section{From Rigid Body To Deforming Body}
\label{sec:rigid-to-deforming}

The rigid body has one orientation \(R(t)\). A Cosserat body assigns an
orientation to each material point:
\[
  R(t)
  \quad\longrightarrow\quad
  R(X,t).
\]
The rigid-body angular velocity
\[
  R^T\dot R
\]
becomes a time connection, while the spatial derivatives
\[
  R^T\partial_iR
\]
become material connection components. This is the first bridge from rigid body
mechanics to gauge theory of deforming bodies.

\section{What To Take Forward}
\label{sec:rigid-take-forward}

\begin{itemize}
\item \(\SO(3)\) is a curved configuration space; Euler angles are useful local
coordinates but not the system itself.
\item Body angular momentum \(M\) and spatial angular momentum \(J\) encode the
same vector in different frames. Confusing them hides the reduction.
\item The free rigid body is integrable because energy and angular momentum
constraints reduce the motion to closed curves and tori.
\item Liouville tori arise from commuting Hamiltonian flows, not from a lucky
choice of variables. Action-angle coordinates are coordinates adapted to those
tori.
\item Replacing one global attitude \(R(t)\) by a field \(R(X,t)\) is the
conceptual step from rigid bodies to gauge descriptions of deformable bodies.
\end{itemize}

\section{Associated Demos}
\label{sec:rigid-associated-demos}

The script \path{demos/python/rigid_body_euler_top.py} integrates Euler's
equations and plots invariant drift. The Mathematica script
\path{demos/mathematica/RigidBodyEulerTop.wl} gives a symbolic and numerical
companion. A useful check is to compare the numerical drift in \(H\) and \(C\)
for several time steps and initial conditions near the intermediate axis.
